\section{Open problems: what a common dynamics would have to do}\label{sec:open}

Almost all open points share one gap: there is no common action, varied in all fields, from which the quantities of the framework
follow.  Up to now they are set.  This section states what such a dynamics would have to do, lists the open problems in that light,
and names the concrete mathematical questions that could be attacked independently of the physical reading.

\subsection{Requirements}

The audited results constrain a dynamics fairly precisely.
\begin{enumerate}
\item \emph{One action for all fields.}  Coframe, connection and matter must be independent fields that act back.  A geometry assembled from
probe fields alone becomes flat (layer v4.9.0-r69).
\item \emph{Matter and tachyon with the same light cone.}  Both need the same kinetic term; the quarter turn \(z\mapsto iz\) must be a symmetry
of the action.
\item \emph{The compensator as a coupling}, not as a projector applied afterwards.
\item \emph{\(\Lambda\) as an integration constant}, as in unimodular gravity~\cite{henneaux} or vacuum-energy sequestering~\cite{kaloper}.
State, frame and normalization never fix the absolute vacuum term (Section~\ref{sec:op5}).
\item \emph{A boundary term that counts one \(U(1)\).}  The inner unit \(J=L_{e_7}\) selects one \(U(1)\) in \(SU(2)\); this might ground the counting
rule behind \(\beta_{\rm sym}\) (conjecture).
\item \emph{Triality as a symmetry of the kinetic term}, broken by chirality.
\item \emph{A clock from the phase difference} between matter and tachyon; a common energy offset is not observable.
\end{enumerate}
A first test of this outline (H48) shows a limit.  In the simplest reduction \(\beta_{\rm sym}\) does not follow from the dynamics, because a
frozen dimensionful value always carries the expansion history.  A dimensionless structural value has to come from a dimensionless ratio
--- two countings (Section~\ref{sec:beta}), or an angle per event.

\subsection{The open problems}

\begin{description}
\item[OP1, the frame offset.]  Derive from the frames why the tachyon frame counts only \(j=\tfrac12\) and the matter frame all spins with weight
two, or find another origin of \(\beta_{\rm sym}\).  Next step: derive or exclude the U(1) selection from the inner unit \(J\).
\item[OP2, a Dirac operator] \(D_-=L_0|_{H^+_-}\) on the antisymmetric sector.  Existence and self-adjointness had been marked as proved on
the basis of the Riemann paper and are open again; the last step, reading the offset from the spectrum, is OP1.
\item[OP3, admissible shell passages.]  Formulated as a local closure certificate and read as compatibility of adjacent cells; for the active
closure phase a local persistence correction is still required.
\item[OP4, the mixing angles.]  Several constructions fit the data, but as fits (Section~\ref{sec:readouts}).  In its old form OP4 compared
values from different frames and was ill-posed; what remains is OP1 together with the pending predictions H42 and H46.
\item[OP5, the vacuum balance.]  Closed locally by a detector threshold, open in exact form.  In the comparison model: decide the possible
further hidden conditions on 7- and 8-jets, by higher precision or by a reformulation that handles the multiplicities of the symbol module
on the characteristic set; then polar spinor variables and the Fierz reconstruction, a constraint algorithm, a coupled coframe dynamics, and
only then a local solution.
\end{description}

Three closure questions belong to the guard list.  The residual bound \(\beta_4=\tfrac1{256}\) of \(C_4\) is a final bound and may not be spent.
For \(C_5\) the multiplicative collar recursion provably has no finite zero; outside it, two of three candidate forms are excluded and the third,
a witness outside the iteration, is open.  The global closure \(C_6/T_2\) is not claimed.  Further open branches are the clock window
\(\Delta_\gamma<\tfrac15\) (derived, measured \(0.1925\) in the matter frame), the dynamics of pair creation (rate, unit of the clock, coupling),
and a reason for three generations beyond the octonions (Section~\ref{sec:octonions}).

\subsection{Questions that stand on their own}

Some of these problems are mathematics that can be studied without accepting any physical reading.
\begin{enumerate}
\item \emph{Formal integrability of the comparison model.}  Is the extended system of Section~\ref{sec:op5} formally integrable, and how many
further conditions appear on 7- and 8-jets?  This is a well-posed question in the formal theory of partial differential equations, with
the data in the repository.
\item \emph{The compensator matrix.}  Find a realization of Definition~\ref{def:cm} on a 16-dimensional carrier, for instance on one generation
of Section~\ref{sec:generation}, that satisfies the hypotheses of Proposition~\ref{prop:spectral}, or show that none exists.
\item \emph{Chirality and anomalies.}  Is the selection of Proposition~\ref{prop:anomaly} known in the literature on
\(\mathbb C\otimes\HH\otimes\OO\), and does it survive other choices of the weak \(SU(2)\)?
\item \emph{Counting rules.}  Is there a horizon counting, in the SU(2) treatment, whose minimal and full versions differ by the factor
\(1.8647\), or a natural reason for the U(1) pair?
\end{enumerate}

\medskip\noindent\emph{Status.}  Everything in this section is open.  The requirements are conclusions from the audited results, not a theory;
they say what a dynamics would have to do, not that one exists.
